% GENERATED FILE. Edit canonical lessons, then run npm run generate:books.
% canonical-source-hash: fnv1a64:617b544412c4fce2
% canonical-lessons: FA-C74-jurab, FA-C74-jib, FA-C74-angoshtar, FA-C74-kolah, FA-C74-kafsh

\chapter{Sock, Pocket, Ring, Hat, Shoe}
\label{ch:fa-sock-pocket-ring-hat-shoe}
\hlchaptermodality{74}

\begin{chapteropening}
\textbf{By the end of this chapter:} \emph{I can say a sock, a pocket, a ring, a hat and a shoe.}

Complete the last lesson of chapter 74: I can say a sock, a pocket, a ring, a hat and a shoe.
\end{chapteropening}

\section[\texorpdfstring{\fa{جوراب} (jurâb)}{جوراب (jurâb)}]{\fa{جوراب} (jurâb) — a sock}

\label{lesson:FA-C74-jurab}



\begin{quote}
Before the new one: say the Persian for a shirt, then the Persian for trousers.
\end{quote}

\subsection*{You'll want to know: \fa{جوراب}}
\textbf{\fa{جوراب}} — \emph{jurâb} — \textquotedblleft{}a sock\textquotedblright{}.

\begin{cousinweb}
\emph{jurâb} has the shape of Arabic \emph{jawrab}, \textquotedblleft{}sock\textquotedblright{}. Lane records the Arabic word as taken from an older Persian \emph{gōrab}, so the word went into Arabic and came back.
\end{cousinweb}

\begin{grammarlens}[title={What you've built}]
One more word for this chapter.
\end{grammarlens}

\subsection*{Your turn}
\begin{itemize}
\raggedright
  \item \emph{Say it:} \emph{jurâb}
  \item \emph{Say it:} \emph{jurâb}, once more
  \item \emph{Say it:} say \emph{jurâb}
  \item \emph{Recall:} say the Persian for a shirt, then the Persian for trousers, then say \emph{jurâb} again
\end{itemize}

\begin{tcolorbox}[breakable,colback=teal!4,colframe=teal!35!black,title={Before you move on}]
What does \fa{جوراب} mean? (\textquotedblleft{}A sock\textquotedblright{}.) Say it once more.
\end{tcolorbox}
\section[\texorpdfstring{\fa{جیب} (jib)}{جیب (jib)}]{\fa{جیب} (jib) — a pocket}

\label{lesson:FA-C74-jib}



\begin{quote}
Before the new one: say the Persian for trousers, then the Persian for a sock.
\end{quote}

\subsection*{You'll want to know: \fa{جیب}}
\textbf{\fa{جیب}} — \emph{jib} — \textquotedblleft{}a pocket\textquotedblright{}.

\begin{cousinweb}
\emph{jib} is borrowed from Arabic \emph{jayb}, \textquotedblleft{}opening of a garment, pocket\textquotedblright{}. Its Arabic root is \emph{j-y-b} (Lane).
\end{cousinweb}

\begin{grammarlens}[title={What you've built}]
One more word for this chapter.
\end{grammarlens}

\subsection*{Your turn}
\begin{itemize}
\raggedright
  \item \emph{Say it:} \emph{jib}
  \item \emph{Say it:} \emph{jib}, once more
  \item \emph{Say it:} say \emph{jib}
  \item \emph{Recall:} say the Persian for trousers, then the Persian for a sock, then say \emph{jib} again
\end{itemize}

\begin{tcolorbox}[breakable,colback=teal!4,colframe=teal!35!black,title={Before you move on}]
What does \fa{جیب} mean? (\textquotedblleft{}A pocket\textquotedblright{}.) Say it once more.
\end{tcolorbox}
\section[\texorpdfstring{\fa{انگشتر} (angoshtar)}{انگشتر (angoshtar)}]{\fa{انگشتر} (angoshtar) — a ring}

\label{lesson:FA-C74-angoshtar}



\begin{quote}
Before the new one: say the Persian for a sock, then the Persian for a pocket.
\end{quote}

\subsection*{You'll want to know: \fa{انگشتر}}
\textbf{\fa{انگشتر}} — \emph{angoshtar} — \textquotedblleft{}a ring\textquotedblright{}.

\begin{cousinweb}
\emph{angoshtar} is built on \emph{angosht}, \textquotedblleft{}finger\textquotedblright{}, a native word usually traced to Middle Persian \emph{angust}: the ring is the finger's piece.
\end{cousinweb}

\begin{grammarlens}[title={What you've built}]
One more word for this chapter.
\end{grammarlens}

\subsection*{Your turn}
\begin{itemize}
\raggedright
  \item \emph{Say it:} \emph{angoshtar}
  \item \emph{Say it:} \emph{angoshtar}, once more
  \item \emph{Say it:} say \emph{angoshtar}
  \item \emph{Recall:} say the Persian for a sock, then the Persian for a pocket, then say \emph{angoshtar} again
\end{itemize}

\begin{tcolorbox}[breakable,colback=teal!4,colframe=teal!35!black,title={Before you move on}]
What does \fa{انگشتر} mean? (\textquotedblleft{}A ring\textquotedblright{}.) Say it once more.
\end{tcolorbox}
\section[\texorpdfstring{\fa{کلاه} (kolâh)}{کلاه (kolâh)}]{\fa{کلاه} (kolâh) — a hat}

\label{lesson:FA-C74-kolah}



\begin{quote}
Before the new one: say the Persian for a pocket, then the Persian for a ring.
\end{quote}

\subsection*{You'll want to know: \fa{کلاه}}
\textbf{\fa{کلاه}} — \emph{kolâh} — \textquotedblleft{}a hat\textquotedblright{}.

\begin{cousinweb}
\emph{kolâh} is a native Persian word, inherited rather than borrowed: it continues Middle Persian \emph{kulāf}, \textquotedblleft{}cap\textquotedblright{} (EtymDB).
\end{cousinweb}

\begin{grammarlens}[title={What you've built}]
One more word for this chapter.
\end{grammarlens}

\subsection*{Your turn}
\begin{itemize}
\raggedright
  \item \emph{Say it:} \emph{kolâh}
  \item \emph{Say it:} \emph{kolâh}, once more
  \item \emph{Say it:} say \emph{kolâh}
  \item \emph{Recall:} say the Persian for a pocket, then the Persian for a ring, then say \emph{kolâh} again
\end{itemize}

\begin{tcolorbox}[breakable,colback=teal!4,colframe=teal!35!black,title={Before you move on}]
What does \fa{کلاه} mean? (\textquotedblleft{}A hat\textquotedblright{}.) Say it once more.
\end{tcolorbox}
\section[\texorpdfstring{\fa{کفش} (kafsh)}{کفش (kafsh)}]{\fa{کفش} (kafsh) — a shoe}

\label{lesson:FA-C74-kafsh}



\begin{quote}
Before the new one: say the Persian for a ring, then the Persian for a hat.
\end{quote}

\subsection*{You'll want to know: \fa{کفش}}
\textbf{\fa{کفش}} — \emph{kafsh} — \textquotedblleft{}a shoe\textquotedblright{}.

\begin{cousinweb}
\emph{kafsh} is a native Persian word, inherited rather than borrowed: it continues Middle Persian \emph{kafš}, \textquotedblleft{}shoe\textquotedblright{} (EtymDB).
\end{cousinweb}

\begin{grammarlens}[title={What you've built}]
One more word for this chapter.
\end{grammarlens}

\subsection*{Your turn}
\begin{itemize}
\raggedright
  \item \emph{Say it:} \emph{kafsh}
  \item \emph{Say it:} \emph{kafsh}, once more
  \item \emph{Say it:} say \emph{kafsh}
  \item \emph{Recall:} say the Persian for a ring, then the Persian for a hat, then say \emph{kafsh} again
\end{itemize}

\begin{tcolorbox}[breakable,colback=teal!4,colframe=teal!35!black,title={Before you move on}]
What does \fa{کفش} mean? (\textquotedblleft{}A shoe\textquotedblright{}.) Say it once more.
\end{tcolorbox}
